\section{Architecture Technique}

\subsection{Vue d'Ensemble de l'Architecture}

DesignPro AI est construit sur une architecture web moderne et scalable, combinant les meilleures pratiques du développement full-stack avec des services d'IA de pointe. La figure \ref{fig:architecture} présente l'architecture globale du système.

\begin{figure}[H]
\centering
\begin{tikzpicture}[
    node distance=1.5cm,
    component/.style={rectangle, draw, fill=blue!20, text width=2.5cm, text centered, minimum height=1cm, font=\small},
    service/.style={rectangle, draw, fill=green!20, text width=2.5cm, text centered, minimum height=1cm, font=\small},
    db/.style={cylinder, draw, fill=orange!20, text width=2cm, text centered, minimum height=1cm, font=\small, shape border rotate=90},
    arrow/.style={->, >=stealth, thick}
]
    % Frontend
    \node[component] (ui) {Interface Utilisateur\\(Next.js + React)};
    
    % Backend
    \node[component, below of=ui, yshift=-0.5cm] (api) {API Routes\\(Next.js Server)};
    
    % Services
    \node[service, left of=api, xshift=-3cm] (mistral) {Mistral AI\\(LLM)};
    \node[service, right of=api, xshift=3cm] (hf) {Hugging Face\\(Stable Diffusion)};
    \node[service, below of=mistral] (openai) {OpenAI\\(GPT-4 Vision)};
    
    % Database
    \node[db, below of=api, yshift=-1cm] (db) {Supabase\\PostgreSQL};
    
    % Arrows
    \draw[arrow] (ui) -- (api);
    \draw[arrow] (api) -- (mistral);
    \draw[arrow] (api) -- (hf);
    \draw[arrow] (api) -- (openai);
    \draw[arrow] (api) -- (db);
\end{tikzpicture}
\caption{Architecture technique de DesignPro AI}
\label{fig:architecture}
\end{figure}

\subsection{Stack Technologique}

\subsubsection{Frontend}

\textbf{Next.js 15} : Framework React avec rendu côté serveur (SSR) et génération statique (SSG)

\textbf{Caractéristiques utilisées} :
\begin{itemize}[leftmargin=*]
    \item \textbf{App Router} : Nouvelle architecture de routage basée sur le système de fichiers
    \item \textbf{Server Components} : Composants React rendus côté serveur pour de meilleures performances
    \item \textbf{Server Actions} : Actions serveur pour les mutations de données
    \item \textbf{Streaming} : Rendu progressif pour une meilleure UX
\end{itemize}

\textbf{React 19} : Bibliothèque UI avec les dernières fonctionnalités

\textbf{TailwindCSS} : Framework CSS utilitaire pour un design moderne et responsive

\textbf{Avantages} :
\begin{itemize}[leftmargin=*]
    \item Performance optimale (SSR + SSG)
    \item SEO-friendly
    \item Developer Experience excellente
    \item Hot Module Replacement (HMR)
\end{itemize}

\subsubsection{Backend et Base de Données}

\textbf{Supabase} : Plateforme Backend-as-a-Service (BaaS) open-source

\textbf{Composants utilisés} :
\begin{itemize}[leftmargin=*]
    \item \textbf{PostgreSQL} : Base de données relationnelle pour le stockage structuré
    \item \textbf{Supabase Auth} : Authentification sécurisée (email/password, OAuth)
    \item \textbf{Row Level Security (RLS)} : Sécurité au niveau des lignes
    \item \textbf{Supabase Storage} : Stockage d'objets pour les images générées
    \item \textbf{Realtime} : Synchronisation en temps réel (optionnel)
\end{itemize}

\textbf{Avantages} :
\begin{itemize}[leftmargin=*]
    \item Open-source et auto-hébergeable
    \item Sécurité intégrée (RLS)
    \item Scalabilité automatique
    \item API REST et GraphQL générées automatiquement
\end{itemize}

\subsubsection{Services d'IA}

\textbf{Mistral AI API} :
\begin{itemize}[leftmargin=*]
    \item Endpoint : \texttt{https://api.mistral.ai/v1/chat/completions}
    \item Modèle : \texttt{mistral-large-latest}
    \item Utilisation : Génération de prompts, évaluations, rapports DfX
\end{itemize}

\textbf{Hugging Face Inference API} :
\begin{itemize}[leftmargin=*]
    \item Endpoint : \texttt{https://router.huggingface.co/hf-inference/models/\{model\_id\}}
    \item Modèles : Stable Diffusion 1.5, 2.1, XL, SDXL-Turbo, OpenJourney
    \item Utilisation : Génération d'images text-to-image et sketch-to-image
\end{itemize}

\textbf{OpenAI API} :
\begin{itemize}[leftmargin=*]
    \item Endpoint : \texttt{https://api.openai.com/v1/chat/completions}
    \item Modèle : \texttt{gpt-4-vision-preview}
    \item Utilisation : Analyse visuelle réelle des designs générés
\end{itemize}

\subsection{Schéma de Données}

\subsubsection{Modèle Conceptuel}

La base de données PostgreSQL est structurée autour de 4 tables principales :

\begin{table}[H]
\centering
\caption{Tables de la base de données}
\label{tab:database}
\begin{tabular}{@{}lp{8cm}@{}}
\toprule
\textbf{Table} & \textbf{Description} \\ 
\midrule
\texttt{users} & Informations utilisateurs (via Supabase Auth) \\
\texttt{projects} & Projets de design (nom, domaine, méthodologie) \\
\texttt{project\_states} & États complets des sessions de design (JSON) \\
\texttt{design\_iterations} & Historique des itérations et scores \\
\bottomrule
\end{tabular}
\end{table}

\subsubsection{Table \texttt{project\_states}}

Cette table centrale stocke l'intégralité du contexte d'une session de design :

\begin{lstlisting}[language=SQL, caption=Schéma de la table project\_states]
CREATE TABLE project_states (
  id UUID PRIMARY KEY DEFAULT uuid_generate_v4(),
  user_id UUID REFERENCES auth.users(id),
  project_id UUID REFERENCES projects(id),
  
  -- Phase 1: Prompt
  methodology VARCHAR(50),
  prompt_generated TEXT,
  
  -- Phase 2: Images
  images JSONB,  -- Array d'URLs d'images
  model_used VARCHAR(100),
  
  -- Phase 3: Evaluation
  agent_q_evaluation JSONB,
  
  -- Phase 4: Simulation
  real_simulation JSONB,
  
  -- Metadata
  created_at TIMESTAMP DEFAULT NOW(),
  updated_at TIMESTAMP DEFAULT NOW()
);
\end{lstlisting}

\subsubsection{Row Level Security (RLS)}

Supabase RLS garantit que chaque utilisateur ne peut accéder qu'à ses propres données :

\begin{lstlisting}[language=SQL, caption=Politique RLS]
-- Politique de lecture
CREATE POLICY "Users can read own project states"
ON project_states FOR SELECT
USING (auth.uid() = user_id);

-- Politique d'écriture
CREATE POLICY "Users can insert own project states"
ON project_states FOR INSERT
WITH CHECK (auth.uid() = user_id);

-- Politique de mise à jour
CREATE POLICY "Users can update own project states"
ON project_states FOR UPDATE
USING (auth.uid() = user_id);
\end{lstlisting}

\subsection{Architecture Client-Serveur}

\subsubsection{Flux de Données}

Le flux de données suit le pattern suivant :

\begin{enumerate}[leftmargin=*]
    \item \textbf{Client} : L'utilisateur interagit avec l'interface React
    \item \textbf{Server Action} : Une action serveur Next.js est déclenchée
    \item \textbf{Service Layer} : Le service \texttt{AIService} orchestre les appels API
    \item \textbf{External APIs} : Appels aux APIs Mistral, Hugging Face, OpenAI
    \item \textbf{Database} : Persistance dans Supabase PostgreSQL
    \item \textbf{Response} : Retour au client avec les données
\end{enumerate}

\subsubsection{Gestion de l'État}

\textbf{Côté Client} :
\begin{itemize}[leftmargin=*]
    \item React State (useState, useReducer)
    \item React Context pour l'état global
    \item SWR pour le cache et la synchronisation
\end{itemize}

\textbf{Côté Serveur} :
\begin{itemize}[leftmargin=*]
    \item Supabase comme source de vérité
    \item Singleton pattern pour \texttt{AIService}
    \item Cache en mémoire pour les modèles chargés
\end{itemize}

\subsection{Intégration des APIs}

\subsubsection{Architecture du Service AI}

Le fichier \texttt{lib/ai.ts} implémente un service singleton \texttt{AIService} qui encapsule toute la logique d'IA :

\begin{lstlisting}[language=TypeScript, caption=Structure du service AI]
export class AIService {
  private static instance: AIService;
  
  public static getInstance(): AIService {
    if (!AIService.instance) {
      AIService.instance = new AIService();
    }
    return AIService.instance;
  }
  
  // Phase 1: Prompt Generation
  async generateProfessionalPrompt(...): Promise<string>
  
  // Phase 2: Image Generation
  async generateDesignCandidates(...): Promise<DesignGenerationResult>
  
  // Phase 3: Evaluation
  async evaluateDesignWithAgentQ(...): Promise<AgentQEvaluation>
  
  // Phase 4: Simulation
  async simulateREALAnalysis(...): Promise<REALSimulation>
}
\end{lstlisting}

\subsubsection{Gestion des Erreurs}

Chaque appel API est enveloppé dans un try-catch avec fallback :

\begin{lstlisting}[language=TypeScript, caption=Pattern de gestion d'erreurs]
async generateDesignCandidates(...) {
  try {
    // Tentative avec modèle primaire
    return await this.generateWithStableDiffusion(prompt);
  } catch (error) {
    console.error('Erreur génération:', error);
    // Fallback vers images de démonstration
    return {
      images: await this.generateRealisticFallbackImages(prompt),
      source: 'local-fallback-error',
      used_fallback: true,
      fallback_reason: error.message
    };
  }
}
\end{lstlisting}

\subsection{Sécurité et Authentification}

\subsubsection{Authentification Supabase}

\textbf{Méthodes supportées} :
\begin{itemize}[leftmargin=*]
    \item Email + Password
    \item OAuth (Google, GitHub, etc.)
    \item Magic Links (email sans mot de passe)
\end{itemize}

\textbf{Flux d'authentification} :
\begin{enumerate}[leftmargin=*]
    \item L'utilisateur se connecte via l'interface
    \item Supabase Auth génère un JWT
    \item Le JWT est stocké dans un cookie HTTP-only
    \item Chaque requête inclut le JWT pour l'autorisation
    \item RLS vérifie les permissions au niveau de la base de données
\end{enumerate}

\subsubsection{Protection des Clés API}

Les clés API sensibles sont stockées dans des variables d'environnement :

\begin{lstlisting}[language=bash, caption=Fichier .env.local]
# Supabase
NEXT_PUBLIC_SUPABASE_URL=https://xxx.supabase.co
NEXT_PUBLIC_SUPABASE_ANON_KEY=eyJxxx...
SUPABASE_SERVICE_ROLE_KEY=eyJxxx...

# AI Services
MISTRAL_API_KEY=xxx
HF_API_TOKEN=hf_xxx
OPENAI_API_KEY=sk-xxx
REPLICATE_API_TOKEN=r8_xxx
\end{lstlisting}

\textbf{Sécurité} :
\begin{itemize}[leftmargin=*]
    \item Les clés ne sont jamais exposées au client
    \item Les appels API se font uniquement côté serveur
    \item Les variables \texttt{NEXT\_PUBLIC\_*} sont les seules accessibles au client
\end{itemize}

\subsection{Déploiement}

\subsubsection{Environnements}

\textbf{Développement} :
\begin{itemize}[leftmargin=*]
    \item Local : \texttt{npm run dev}
    \item Hot reload activé
    \item Logs détaillés
\end{itemize}

\textbf{Production} :
\begin{itemize}[leftmargin=*]
    \item Plateforme : Vercel (recommandé pour Next.js)
    \item Build optimisé : \texttt{npm run build}
    \item Déploiement automatique via Git
\end{itemize}

\subsubsection{Configuration Vercel}

\begin{lstlisting}[language=JSON, caption=vercel.json]
{
  "buildCommand": "npm run build",
  "outputDirectory": ".next",
  "framework": "nextjs",
  "env": {
    "MISTRAL_API_KEY": "@mistral-api-key",
    "HF_API_TOKEN": "@hf-api-token",
    "OPENAI_API_KEY": "@openai-api-key"
  }
}
\end{lstlisting}

\subsection{Performance et Optimisation}

\subsubsection{Optimisations Frontend}

\begin{itemize}[leftmargin=*]
    \item \textbf{Code Splitting} : Chargement lazy des composants lourds
    \item \textbf{Image Optimization} : Next.js Image component pour les images optimisées
    \item \textbf{Caching} : SWR pour le cache des requêtes
    \item \textbf{Compression} : Gzip/Brotli activé sur Vercel
\end{itemize}

\subsubsection{Optimisations Backend}

\begin{itemize}[leftmargin=*]
    \item \textbf{Connection Pooling} : Supabase gère automatiquement le pool de connexions
    \item \textbf{Model Caching} : Les modèles Hugging Face sont mis en cache côté serveur
    \item \textbf{Retry Logic} : Réessais automatiques avec backoff exponentiel
    \item \textbf{Low Memory Mode} : Option pour réduire l'utilisation mémoire
\end{itemize}

\subsection{Monitoring et Logs}

\subsubsection{Logging}

Système de logs structuré avec différents niveaux :

\begin{lstlisting}[language=TypeScript, caption=Système de logging]
console.log('🎨 Lancement Phase 2: Génération...');
console.log('✅ Évaluation Agent Q générée:', score);
console.warn('⚠️ Aucune API d\'analyse visuelle disponible');
console.error('❌ Erreur critique génération Phase 2:', error);
\end{lstlisting}

\subsubsection{Métriques}

\textbf{Métriques collectées} :
\begin{itemize}[leftmargin=*]
    \item Temps de génération par phase
    \item Taux de succès/échec des API
    \item Utilisation des modèles (primaire vs fallback)
    \item Scores utilisateurs (feedback)
\end{itemize}

Ces métriques permettent d'optimiser continuellement le système et d'identifier les points de friction.
